\section{Brownian Compression in the Frozen Fluid Teacher}\label{app:ns-compression}
We estimate how much endpoint variation the finite Brownian features omit, independently of training a neural map. Fix validation trajectory 160 at origin 45 and its five-field history. For each feature count $K\in\{1,4,16\}$, we sample eight coarse paths and four independent finer Brownian bridge realizations conditional on each coarse path. The frozen SI teacher integrates its drift using 64 Euler steps with exact martingale increments. Cached coarse observations are copied into independent bridge forks; new fine values in a fork are not fed back into its parent. Thus replicas share the prescribed coarse features and differ only in unobserved path details.

Let $X_{gr}$ be an endpoint for coarse group $g$ and replica $r$, and let $\bar X_g$ be its group mean. Compute coordinatewise sample variances and then average over the state coordinates:
\begin{align}
 \widehat W&=\frac1G\sum_{g=1}^G
 \left\langle\operatorname{Var}_{r}(X_{gr})\right\rangle,\\
 \widehat B&=\left\langle\operatorname{Var}_{g}(\bar X_g)\right\rangle,
 \qquad G=8,\quad R=4.
\end{align}
All sample variances use their unbiased denominators. The first statistic estimates the conditional-variance term in Proposition~\ref{prop:compression} for this fixed history and numerical teacher. Finite replication inflates the variance of group means by $\widehat W/R$, so we estimate explained variance as $\widehat B-\widehat W/R$ and total variance as $\widehat B+(1-1/R)\widehat W$.

\begin{figure}[t]
\centering
\pendingbox{mean within-group variance per coarse Brownian path (points) and mean $\pm$ one standard error across the eight groups (bars), for $K=1,4,16$.}
\caption{Fixed-history teacher diagnostic, not forecast test performance. All coordinates are normalized. More path features expose Brownian variation otherwise hidden from a same-path predictor.}
\label{fig:ns-compression}
\end{figure}

\pending{estimated omitted variances $\widehat W$ and their standard errors for $K=1,4,16$; omitted fractions of estimated total endpoint variance; reduction factor from $K=4$ to $K=16$ (Figure~\ref{fig:ns-compression}).}

This diagnostic uses only one history and a small nested Monte Carlo sample. True total conditional forecast variance is independent of $K$, so variation of its estimates across $K$ indicates Monte Carlo uncertainty in the fractions. The standard errors do not quantify variation across histories, and grid cells are not treated as independent replicates. The teacher itself is learned and numerically integrated, and its 64-step diagnostic differs from the 32-step supervision and the original SI sampler's Euler--Maruyama increments. The estimates cannot be subtracted directly from aggregate student validation losses.

Finally, this is an information limit for squared-error prediction under the specified Brownian coupling. It is neither a lower bound on CRPS nor a guarantee that increasing $K$ improves trained-map forecasts. The finite-feature PSD objective also need not have the same optimum as conditional-mean endpoint regression. \pending{fraction of the teacher's standard deviation retained by the estimated $K=4$ conditional-mean predictor, and whether compression alone accounts for any loss of spread in the trained students.}
